\documentclass[12pt,a4paper]{report}

% ── Encodage et langue (XeLaTeX) ────────────────────────────────
\usepackage{fontspec}
\usepackage{polyglossia}
\setmainlanguage{french}
\setmainfont{DejaVu Serif}
\setmonofont{DejaVu Sans Mono}

% ── Mise en page ────────────────────────────────────────────────
\usepackage[top=2.5cm,bottom=2.5cm,left=3cm,right=2.5cm]{geometry}
\usepackage{setspace}
\onehalfspacing

% ── Couleurs ────────────────────────────────────────────────────
\usepackage{xcolor}
\definecolor{accentColor}{HTML}{D4A017}
\definecolor{darkGray}{HTML}{2D2D2D}
\definecolor{codeGray}{HTML}{F5F5F5}
\definecolor{linkColor}{HTML}{1A56A0}

% ── Hyperliens ──────────────────────────────────────────────────
\usepackage[colorlinks=true,linkcolor=linkColor,urlcolor=linkColor]{hyperref}

% ── En-têtes et pieds de page ───────────────────────────────────
\usepackage{fancyhdr}
\setlength{\headheight}{15pt}
\pagestyle{fancy}
\fancyhf{}
\fancyhead[L]{\small\textit{Système de Facturation PHP -- UPC 2025-2026}}
\fancyhead[R]{\small\thepage}
\fancyfoot[C]{\small\textit{Travaux Pratiques -- Programmation Web PHP}}
\renewcommand{\headrulewidth}{0.4pt}
\renewcommand{\footrulewidth}{0.4pt}

% ── Code source ─────────────────────────────────────────────────
\usepackage{listings}
\lstset{
  basicstyle=\footnotesize\ttfamily,
  backgroundcolor=\color{codeGray},
  frame=single,
  framesep=4pt,
  rulecolor=\color{gray!40},
  keywordstyle=\color{blue!70}\bfseries,
  commentstyle=\color{green!50!black}\itshape,
  stringstyle=\color{red!60!black},
  breaklines=true,
  tabsize=4,
  showstringspaces=false,
  numbers=left,
  numberstyle=\tiny\color{gray},
  numbersep=8pt,
}

% ── Tableaux ────────────────────────────────────────────────────
\usepackage{booktabs}
\usepackage{tabularx}
\usepackage{array}
\newcolumntype{L}[1]{>{\raggedright\arraybackslash}p{#1}}

% ── Figures ─────────────────────────────────────────────────────
\usepackage{graphicx}
\usepackage{float}
\usepackage{tikz}
\usetikzlibrary{shapes.geometric,arrows.meta,positioning}

% ── Listes ──────────────────────────────────────────────────────
\usepackage{enumitem}

% ── Boîtes colorées ─────────────────────────────────────────────
\usepackage{tcolorbox}
\tcbuselibrary{skins,breakable}
\newtcolorbox{notebox}{
  colback=accentColor!10,colframe=accentColor,
  fonttitle=\bfseries, title=Remarque,
  breakable
}
\newtcolorbox{warningbox}{
  colback=red!5,colframe=red!60!black,
  fonttitle=\bfseries, title=Attention,
  breakable
}

% ───────────────────────────────────────────────────────────────
\begin{document}

% ╔═══════════════════════════════╗
% ║       PAGE DE GARDE          ║
% ╚═══════════════════════════════╝
\begin{titlepage}
  \centering
  \vspace*{2cm}

  {\Large\bfseries Université Protestante au Congo}\\[0.4cm]
  {\large Faculté de Sciences Informatiques}\\[0.2cm]
  {\large Département Informatique -- L2 FASI}\\[2cm]

  \rule{\linewidth}{1.5pt}\\[0.5cm]
  {\huge\bfseries\color{darkGray} Système de Facturation}\\[0.4cm]
  {\Large avec Lecture de Codes-Barres}\\[0.3cm]
  \rule{\linewidth}{1.5pt}\\[1.5cm]

  {\large\textit{Travaux Pratiques de Programmation Web en PHP}}\\[0.3cm]
  {\normalsize Programmation Procédurale -- Gestion de Fichiers -- Contrôle d'Accès}\\[2cm]

  \begin{tabular}{ll}
    \textbf{Auteurs :}        & [Étudiant 1 -- Nom Prénom]          \\[0.2cm]
                              & [Étudiant 2 -- Nom Prénom]          \\[0.2cm]
                              & [Étudiant 3 -- Nom Prénom]          \\[0.6cm]
    \textbf{Niveau :}         & Licence 2 -- FASI                   \\[0.2cm]
    \textbf{Matière :}        & Programmation Web en PHP            \\[0.2cm]
    \textbf{Enseignant :}     & [Nom de l'enseignant]               \\[0.2cm]
    \textbf{Année :}          & 2025 -- 2026                        \\[0.2cm]
    \textbf{Date de remise :} & 24 Avril 2026                       \\
  \end{tabular}

  \vfill
  {\small Code source disponible sur GitHub :
  \url{https://github.com/[votre-repo]/facturation}}
\end{titlepage}

\tableofcontents
\newpage

% ╔═══════════════════════════════╗
% ║     CHAPITRE 1 : INTRO       ║
% ╚═══════════════════════════════╝
\chapter{Introduction}

\section{Présentation du projet}

Ce projet consiste en la conception et le développement d'un système de caisse
informatisé pour un supermarché de taille modeste. L'objectif principal est de
remplacer la saisie manuelle des articles -- source d'erreurs et de lenteur -- par
un système de lecture automatique de codes-barres via la caméra d'un téléphone ou
d'un ordinateur.

Le système est entièrement développé en \textbf{PHP procédural}, sans recours à
aucun système de gestion de base de données. La persistance des données repose
exclusivement sur des fichiers au format JSON, ce qui garantit une portabilité
maximale et facilite le déploiement sur des serveurs d'hébergement simples.

\section{Contexte}

Le gérant du supermarché souhaitait moderniser son processus de vente afin de :
\begin{itemize}
  \item Réduire les erreurs de saisie commises par les caissiers ;
  \item Accélérer les transactions à la caisse ;
  \item Disposer d'un historique structuré des ventes ;
  \item Contrôler en temps réel les niveaux de stock.
\end{itemize}

La contrainte budgétaire imposait l'absence de base de données externe, rendant
obligatoire l'utilisation de fichiers comme mécanisme de persistance.

\section{Objectifs pédagogiques}

Ce travail pratique nous a permis de mettre en oeuvre les compétences suivantes :
\begin{enumerate}
  \item Concevoir une application PHP selon le paradigme \textbf{procédural} ;
  \item Utiliser les \textbf{fichiers JSON} comme base de données légère ;
  \item Implémenter un \textbf{contrôle d'accès basé sur les rôles} (RBAC) ;
  \item Intégrer la bibliothèque JavaScript \textbf{QuaggaJS} pour la lecture de codes-barres ;
  \item Appliquer les bonnes pratiques de \textbf{validation et d'assainissement} des données ;
  \item Rédiger un rapport technique structuré en \LaTeX{}.
\end{enumerate}

% ╔═══════════════════════════════╗
% ║  CHAPITRE 2 : STRUCTURE      ║
% ╚═══════════════════════════════╝
\chapter{Structure du Projet}

\section{Arborescence commentée}

L'arborescence du projet suit une organisation modulaire qui sépare clairement les
responsabilités : configuration, authentification, modules fonctionnels, données,
fonctions réutilisables et ressources statiques.

\begin{lstlisting}[caption={Arborescence du projet}]
facturation/
|-- index.php                   Tableau de bord principal
|-- README.txt                  Instructions de deploiement
|-- config/
|   `-- config.php              Constantes globales : TVA, chemins, roles
|-- auth/
|   |-- login.php               Formulaire d'authentification
|   |-- logout.php              Destruction de session et redirection
|   `-- session.php             Gestion des sessions et controle d'acces
|-- modules/
|   |-- produits/
|   |   |-- enregistrer.php     Enregistrement / MAJ d'un produit
|   |   |-- lire.php            API AJAX : lecture produit par code-barres
|   |   `-- liste.php           Catalogue avec gestion du stock
|   |-- facturation/
|   |   |-- nouvelle-facture.php  Interface de caisse
|   |   |-- calcul.php            API AJAX : calcul des totaux
|   |   `-- afficher-facture.php  Affichage et impression
|   `-- admin/
|       |-- gestion-comptes.php   Liste des comptes utilisateurs
|       |-- ajouter-compte.php    Creation d'un nouveau compte
|       `-- supprimer-compte.php  Desactivation avec confirmation
|-- data/
|   |-- produits.json           Catalogue produits (persistance)
|   |-- factures.json           Historique des factures (persistance)
|   `-- utilisateurs.json       Comptes utilisateurs avec hash bcrypt
|-- includes/
|   |-- header.php              En-tete HTML commun
|   |-- footer.php              Pied de page HTML commun
|   |-- fonctions-produits.php  CRUD produits, validation, stock
|   |-- fonctions-factures.php  Creation, calculs, lecture des factures
|   `-- fonctions-auth.php      Authentification, CRUD utilisateurs
|-- assets/
|   |-- css/style.css           Feuille de style principale (theme sombre)
|   `-- js/scanner.js           Encapsulation QuaggaJS
`-- rapports/
    |-- rapport-journalier.php  Rapport des ventes du jour
    |-- rapport-mensuel.php     Rapport des ventes par mois
    `-- rapport.tex             Rapport technique (ce document)
\end{lstlisting}

\section{Description détaillée des fichiers}

\subsection{Fichier \texttt{config/config.php}}
Ce fichier centralise toutes les constantes de l'application : taux de TVA (18\,\%),
chemins absolus vers les fichiers de données, noms des rôles, hiérarchie des rôles
et préfixe des identifiants de facture. Son inclusion en premier dans chaque script
garantit la cohérence des paramètres à travers toute l'application.

\subsection{Fichier \texttt{auth/session.php}}
Il gère le cycle de vie de la session PHP. Ses fonctions clés sont :
\begin{itemize}
  \item \texttt{verifier\_connecte()} : redirige vers la page de connexion si aucune session n'est active ;
  \item \texttt{verifier\_role(\$role)} : vérifie que l'utilisateur possède le niveau de permission suffisant ;
  \item \texttt{a\_role(\$role)} : retourne un booléen, utilisé dans les vues pour afficher conditionnellement les menus ;
  \item Les fonctions \texttt{chemin\_base()}, \texttt{chemin\_auth()}, \texttt{chemin\_module()} calculent les chemins relatifs HTML depuis n'importe quel sous-répertoire.
\end{itemize}

\subsection{Fichier \texttt{includes/fonctions-produits.php}}
Contient la logique métier liée aux produits :
\texttt{charger\_produits()}, \texttt{sauvegarder\_produits()}, \texttt{trouver\_produit()},
\texttt{enregistrer\_produit()}, \texttt{decrementer\_stock()}, \texttt{modifier\_stock()},
\texttt{valider\_produit()} et \texttt{formater\_prix()}.

\subsection{Fichier \texttt{includes/fonctions-factures.php}}
Gère la création et la lecture des factures :
\texttt{charger\_factures()}, \texttt{sauvegarder\_factures()}, \texttt{generer\_id\_facture()},
\texttt{calculer\_totaux()}, \texttt{creer\_facture()}, \texttt{trouver\_facture()},
\texttt{factures\_du\_jour()} et \texttt{factures\_du\_mois()}.

\subsection{Fichier \texttt{modules/facturation/calcul.php}}
Point de terminaison AJAX appelé depuis l'interface de caisse pour valider les articles
et calculer les totaux côté serveur avant l'enregistrement définitif. Ce double
contrôle serveur évite toute manipulation côté client.

\subsection{Fichier \texttt{assets/js/scanner.js}}
Module JavaScript qui encapsule QuaggaJS. Il expose \texttt{initScanner(containerId,
onDetected, statusId)} et \texttt{stopScanner()}. Pour éviter les faux positifs,
un code-barres n'est confirmé qu'après \textbf{trois lectures identiques consécutives}.

% ╔═══════════════════════════════╗
% ║  CHAPITRE 3 : DONNÉES        ║
% ╚═══════════════════════════════╝
\chapter{Structures de Données}

\section{Justification du format JSON}

Le format JSON a été retenu pour plusieurs raisons :
\begin{enumerate}
  \item \textbf{Lisibilité humaine} : les fichiers peuvent être inspectés et modifiés manuellement ;
  \item \textbf{Support natif en PHP} : \texttt{json\_encode()} et \texttt{json\_decode()} sans extension externe ;
  \item \textbf{Portabilité} : aucune dépendance à un serveur de base de données ;
  \item \textbf{Structures hiérarchiques} : les articles d'une facture sont imbriqués directement, sans jointures.
\end{enumerate}

\begin{notebox}
  Les fichiers JSON sont écrits avec \texttt{JSON\_PRETTY\_PRINT} pour la lisibilité
  et \texttt{JSON\_UNESCAPED\_UNICODE} pour préserver les caractères accentués.
\end{notebox}

\section{Structure d'un produit}

\begin{lstlisting}[language=,caption={Structure JSON d'un produit}]
{
  "code_barre":          "3017620422003",
  "nom":                 "Huile de palme 1L",
  "prix_unitaire_ht":    1200,
  "date_expiration":     "2027-12-31",
  "quantite_stock":      50,
  "date_enregistrement": "2026-04-17"
}
\end{lstlisting}

Le champ \texttt{code\_barre} sert de clé primaire implicite.

\section{Structure d'une facture}

\begin{lstlisting}[language=,caption={Structure JSON d'une facture}]
{
  "id_facture": "FAC-20260417-001",
  "date":       "2026-04-17",
  "heure":      "10:35:22",
  "caissier":   "caissier.demo",
  "articles": [
    {
      "code_barre":       "3017620422003",
      "nom":              "Huile de palme 1L",
      "prix_unitaire_ht": 1200,
      "quantite":         2,
      "sous_total_ht":    2400
    }
  ],
  "total_ht":  2400,
  "tva":        432,
  "total_ttc": 2832
}
\end{lstlisting}

L'identifiant est généré au format \texttt{FAC-AAAAMMJJ-NNN}, où \texttt{NNN}
est un compteur journalier sur trois chiffres, garantissant l'unicité sans table de séquence.

\section{Structure d'un utilisateur}

\begin{lstlisting}[language=,caption={Structure JSON d'un utilisateur}]
{
  "identifiant":    "jean.mbeki",
  "mot_de_passe":   "$2y$10$...",
  "role":           "caissier",
  "nom_complet":    "Jean Mbeki",
  "date_creation":  "2026-04-17",
  "actif":          true
}
\end{lstlisting}

Le mot de passe est haché via \texttt{password\_hash(\$mdp, PASSWORD\_DEFAULT)}
(algorithme bcrypt, coût 10). Le champ \texttt{actif} permet une suppression logique.

% ╔═══════════════════════════════╗
% ║  CHAPITRE 4 : MODULES        ║
% ╚═══════════════════════════════╝
\chapter{Description des Modules}

\section{Module d'authentification}

\subsection{Mécanisme de connexion}
La page \texttt{auth/login.php} reçoit les identifiants via POST. La fonction
\texttt{authentifier\_utilisateur()} parcourt le fichier JSON et appelle
\texttt{password\_verify()} pour chaque entrée correspondante. Si l'authentification
réussit, les données de l'utilisateur (sans le hash) sont stockées dans
\texttt{\$\_SESSION['utilisateur']}.

\subsection{Protection des pages}
Chaque script inclut \texttt{auth/session.php} et appelle \texttt{verifier\_role(\$role\_minimum)}.
La hiérarchie numérique permet une comparaison simple :

\begin{lstlisting}[language=PHP,caption={Hiérarchie des rôles}]
define('ROLES_HIERARCHIE', [
    ROLE_CAISSIER   => 1,
    ROLE_MANAGER    => 2,
    ROLE_SUPERADMIN => 3,
]);
\end{lstlisting}

\section{Module d'enregistrement des produits}

\subsection{Flux de traitement}
\begin{enumerate}
  \item Le Manager active la caméra -- QuaggaJS lit le code-barres ;
  \item Le JavaScript redirige vers \texttt{enregistrer.php?code=...} ;
  \item PHP appelle \texttt{trouver\_produit(\$code)} ;
  \item Si inconnu : formulaire vierge ; si connu : formulaire pré-rempli ;
  \item À la soumission : \texttt{valider\_produit()} puis \texttt{enregistrer\_produit()}.
\end{enumerate}

\subsection{Validation des données}
La fonction \texttt{valider\_produit()} vérifie : code-barres non vide, nom non vide,
prix numérique et positif, quantité non négative, date au format \texttt{AAAA-MM-JJ}
via \texttt{checkdate()}.

\section{Module de facturation}

\subsection{Architecture côté client}
L'interface utilise un tableau JavaScript \texttt{articles[]} qui s'accumule au fil
des scans. Chaque ajout déclenche le recalcul des totaux en temps réel, sans rechargement
de page. À la validation, le tableau est sérialisé en JSON et envoyé au serveur.

\subsection{Calcul des montants}

\begin{lstlisting}[language=PHP,caption={Calcul des totaux}]
function calculer_totaux(array $articles): array {
    $total_ht = array_sum(array_column($articles, 'sous_total_ht'));
    $tva      = round($total_ht * TAUX_TVA);
    return [
        'total_ht'  => $total_ht,
        'tva'       => $tva,
        'total_ttc' => $total_ht + $tva,
    ];
}
\end{lstlisting}

\subsection{Format d'affichage d'une facture}

\begin{center}
\begin{tabular}{L{4cm} r r r}
  \toprule
  \textbf{Désignation} & \textbf{Prix HT} & \textbf{Qté} & \textbf{Sous-total HT} \\
  \midrule
  Huile de palme 1L    & 1\,200 CDF       & 2            & 2\,400 CDF             \\
  Savon de Marseille   & 450 CDF          & 5            & 2\,250 CDF             \\
  \midrule
  \multicolumn{3}{r}{Total HT}   & 4\,650 CDF \\
  \multicolumn{3}{r}{TVA (18\%)} & 837 CDF    \\
  \midrule
  \multicolumn{3}{r}{\textbf{Net à payer}} & \textbf{5\,487 CDF} \\
  \bottomrule
\end{tabular}
\end{center}

\subsection{Protection du stock}
Avant l'enregistrement, le serveur re-vérifie le stock de chaque article. Si insuffisant,
la transaction est bloquée. Le stock est décrémenté \emph{après} cette vérification dans
\texttt{decrementer\_stock()}.

\section{Module d'administration des comptes}

\subsection{Tableau des rôles et permissions}

\begin{center}
\begin{tabular}{L{2.5cm} L{9cm}}
  \toprule
  \textbf{Rôle} & \textbf{Permissions} \\
  \midrule
  Caissier  & Lecture codes-barres, création et consultation de factures \\
  Manager   & Idem + enregistrement produits, gestion stock, rapports \\
  Super Admin & Idem + création et désactivation de comptes utilisateurs \\
  \bottomrule
\end{tabular}
\end{center}

\subsection{Sécurité des mots de passe}
La création d'un compte appelle \texttt{password\_hash(\$mdp, PASSWORD\_DEFAULT)}.
En PHP 8, \texttt{PASSWORD\_DEFAULT} correspond à bcrypt (coût 10), résistant aux
attaques par force brute et aux tables arc-en-ciel.

% ╔═══════════════════════════════╗
% ║  CHAPITRE 5 : FLUX           ║
% ╚═══════════════════════════════╝
\chapter{Diagramme de Flux}

\section{Transaction de vente -- vue d'ensemble}

\begin{figure}[H]
\centering
\begin{tikzpicture}[
  node distance=1.1cm,
  startstop/.style={rectangle, rounded corners=6pt, draw=green!60!black,
    fill=green!10, minimum width=3.8cm, minimum height=0.75cm, font=\small\bfseries, align=center},
  process/.style={rectangle, draw=blue!60!black, fill=blue!8,
    minimum width=4cm, minimum height=0.75cm, font=\small, align=center},
  decision/.style={diamond, aspect=2.8, draw=orange!80!black, fill=orange!10,
    font=\scriptsize, inner sep=1pt, align=center},
  arrow/.style={-Stealth, thick}
]

\node[startstop] (start)  {Caissier active\\le scanner};
\node[process,   below=of start]   (scan)   {QuaggaJS lit\\le code-barres};
\node[decision,  below=of scan]    (connu)  {Produit\\connu ?};
\node[process,   right=3cm of connu]  (msg)   {Message :\\enregistrer\\d'abord};
\node[process,   below=of connu]   (affiche){Afficher nom\\et prix HT};
\node[process,   below=of affiche] (qte)    {Saisir la\\quantité};
\node[decision,  below=of qte]     (stock)  {Stock\\suffisant ?};
\node[process,   right=3cm of stock]  (alerte){Alerte\\stock};
\node[process,   below=of stock]   (ajout)  {Ajouter à\\la facture};
\node[decision,  below=of ajout]   (autre)  {Autre\\article ?};
\node[process,   below=of autre]   (valide) {Valider\\la facture};
\node[process,   below=of valide]  (dec)    {Décrémenter\\les stocks};
\node[process,   below=of dec]     (save)   {Sauvegarder\\factures.json};
\node[startstop, below=of save]    (fin)    {Afficher /\\imprimer};

\draw[arrow] (start) -- (scan);
\draw[arrow] (scan) -- (connu);
\draw[arrow] (connu) -- node[left,font=\scriptsize]{Oui} (affiche);
\draw[arrow] (connu) -- node[above,font=\scriptsize]{Non} (msg);
\draw[arrow] (affiche) -- (qte);
\draw[arrow] (qte) -- (stock);
\draw[arrow] (stock) -- node[left,font=\scriptsize]{Oui} (ajout);
\draw[arrow] (stock) -- node[above,font=\scriptsize]{Non} (alerte);
\draw[arrow] (ajout) -- (autre);
\draw[arrow] (autre) -- node[left,font=\scriptsize]{Non} (valide);
\draw[arrow] (autre.east) to[out=0,in=0,looseness=2] node[right,font=\scriptsize]{Oui} (scan.east);
\draw[arrow] (valide) -- (dec);
\draw[arrow] (dec) -- (save);
\draw[arrow] (save) -- (fin);

\end{tikzpicture}
\caption{Flux complet d'une transaction de vente}
\end{figure}

% ╔═══════════════════════════════╗
% ║  CHAPITRE 6 : DIFFICULTÉS    ║
% ╚═══════════════════════════════╝
\chapter{Difficultés Rencontrées et Solutions}

\section{Faux positifs du scanner de codes-barres}

\textbf{Problème :} QuaggaJS déclenchait le callback plusieurs fois pour un même
code-barres en quelques millisecondes, causant des doublons dans la facture.

\textbf{Solution :} Un tampon de lectures (buffer) est implémenté dans \texttt{scanner.js}.
Le code n'est confirmé qu'après \textbf{trois lectures identiques consécutives}
(\texttt{THRESHOLD = 3}). Le buffer est réinitialisé après chaque confirmation.

\section{Gestion des chemins relatifs multi-niveaux}

\textbf{Problème :} Les inclusions PHP et les liens HTML devaient fonctionner depuis
des fichiers à différents niveaux de profondeur dans l'arborescence.

\textbf{Solution :} Utilisation de \texttt{\_\_DIR\_\_} pour les inclusions PHP
(chemin absolu du fichier courant) et fonctions utilitaires \texttt{chemin\_base()},
\texttt{chemin\_auth()}, \texttt{chemin\_module()} dans \texttt{session.php} pour
les liens HTML.

\section{Accès à la caméra en HTTP}

\textbf{Problème :} Les navigateurs modernes bloquent l'accès à la caméra sur HTTP
non sécurisé, sauf sur \texttt{localhost}.

\textbf{Solution :} En développement, \texttt{localhost} contourne la restriction.
En production, un certificat TLS (HTTPS) est obligatoire. Ce point est documenté
dans le \texttt{README.txt}.

\section{Génération d'identifiants de factures uniques}

\textbf{Problème :} Sans table de séquences, le compteur journalier \texttt{NNN}
dans \texttt{FAC-AAAAMMJJ-NNN} devait être calculé dynamiquement.

\textbf{Solution :} La fonction \texttt{generer\_id\_facture()} parcourt toutes
les factures existantes pour trouver le numéro maximum de la journée courante,
puis l'incrémente. Fiable pour des volumes inférieurs à 999 factures par jour.

\section{Concurrence d'écriture}

\textbf{Problème :} Deux caissiers validant une facture simultanément risquent de
corrompre le fichier JSON par écritures concurrentes.

\textbf{Solution appliquée :} Pour ce projet à usage modéré, la probabilité de
collision est faible. Une amélioration future consistera à utiliser le drapeau
\texttt{LOCK\_EX} dans \texttt{file\_put\_contents()}, ou à migrer vers SQLite.

% ╔═══════════════════════════════╗
% ║  CHAPITRE 7 : CONCLUSION     ║
% ╚═══════════════════════════════╝
\chapter{Conclusion et Pistes d'Amélioration}

\section{Bilan du projet}

Ce projet nous a permis de mettre en pratique l'ensemble des concepts du cours de
Programmation Web PHP : architecture procédurale, gestion de fichiers, sessions,
sécurité (RBAC, hachage bcrypt), intégration d'une bibliothèque JavaScript externe
et validation rigoureuse des données.

Le résultat est un système de caisse opérationnel avec :
\begin{itemize}
  \item Lecture de codes-barres via caméra (QuaggaJS) ;
  \item Gestion complète du catalogue produits ;
  \item Facturation avec calcul de TVA à 18\,\% ;
  \item Contrôle d'accès à trois niveaux (Caissier / Manager / Super Admin) ;
  \item Rapports journaliers et mensuels ;
  \item Impression des factures.
\end{itemize}

\section{Pistes d'amélioration}

\subsection{Court terme}
\begin{itemize}
  \item Ajouter le verrouillage de fichier (\texttt{LOCK\_EX}) pour les environnements concurrents ;
  \item Export des rapports en PDF côté serveur (FPDF ou DomPDF) ;
  \item Réinitialisation du mot de passe par courrier électronique.
\end{itemize}

\subsection{Moyen terme}
\begin{itemize}
  \item Migration vers SQLite pour de meilleures performances avec de grands volumes ;
  \item Module de gestion des fournisseurs et des commandes de réapprovisionnement ;
  \item Tableau de bord analytique avec graphiques (Chart.js).
\end{itemize}

\subsection{Long terme}
\begin{itemize}
  \item Réécriture en PHP orienté objet avec un framework MVC (Laravel ou Symfony) ;
  \item Application mobile native pour la lecture de codes-barres ;
  \item Connexion à une imprimante thermique de tickets de caisse.
\end{itemize}

\vspace{2cm}
\begin{center}
  \rule{0.5\linewidth}{0.5pt}\\[0.5cm]
  {\itshape << La réussite se forge par l'effort et la persévérance. >>}\\[0.3cm]
  {\small Université Protestante au Congo -- Faculté de Sciences Informatiques -- 2025-2026}
\end{center}

\end{document}
